% =====================================================================
%  GMN2 §3.1  从膜构造到 Hitchin 系统
%  原文：wkb-paper.tex 1374–2190 行；印刷 p.19–31；PDF p.20–32
%  编号公式 (3.1)–(3.40) 共 40 条（含原文一处 3 行 align 拆成 3 个 equation）；脚注 6 条；插图 4 幅
%  本块内 \eqref/\ref 指向本文件内的 label；跨块引用硬编码原文编号并加 % TODO
%  插图文件名：string / m2-branch / bps-disc / bps-annulus（图 1–4）
% =====================================================================

\section{从膜构造到 Hitchin 系统} \label{sec:physics}

本节我们描述一类 $d=4$、$\CN=2$ 场论，其超凯勒流形 $\CM$ 是 Hitchin 方程之解的模空间。

在第 \ref{sec:scft} 节中，我们以六维 $(2,0)$ 理论在黎曼面 $C$ 上约化的语言来讨论这些理论——$C$ 的若干点上还带有实余维 2 的缺陷——并解释为什么把它们进一步紧化到圆上就会给出 Hitchin 方程。在第 3.2 节% TODO 合稿改 \ref{subsec:WittenConstruction}
中，我们通过与 \cite{Witten:1997sc} 所述 D4/D6/NS5-膜构造对接，给出我们所考虑的一大类理论的纯四维定义。最后，在第 3.3 节% TODO 合稿改 \ref{sec:iib}
中，我们简要回顾另一种会给出同一批理论的 Type IIB 弦论构造，其中涉及一个以 $C$ 为 ADE 奇点曲线的 Calabi–Yau 三维流形。

\subsection{紧化 $(2,0)$ 理论} \label{sec:scft}

从著名的六维 $(2,0)$ 超共形场论出发，在适当的黎曼面 $C$ 上紧化，就能得到极大一类 $d=4$、$\CN=2$ 场论。正如我们将在下面第 \ref{subsubsec:20AndCoulomb}–\ref{subsubsec:KineticTerms} 节中所描述的，这一构造还自然地把这些理论中的 Seiberg–Witten 曲线实现为 $T^*C$ 内 $C$ 的分歧覆盖，同时给出直接把握 BPS 谱的办法，从而推广了 \cite{Klemm:1996bj} 中描述的一幅图像。

本文的主要兴趣在于：把这些四维理论进一步紧化到圆上会发生什么。也就是说，我们考虑 $\IR^{1,2} \times S^1 \times C$ 上的 $(2,0)$ 理论，其中 $S^1$ 的半径为 $R$。于是低能下我们得到一个三维 sigma 模型，其靶空间 $\CM$ 是超凯勒流形。

$\CM$ 是什么？为回答这个问题，我们把同一个理论换一种方式来描述，即颠倒紧化的顺序。若先把六维 $(2,0)$ 理论紧化在 $S^1$ 上，低能物理由五维超对称 Yang–Mills 理论描述。随后我们可以把这个五维理论进一步紧化在 $C$ 上。从这个观点看，三维超庞加莱不变的真空所成的模空间，就是 $C$ 上规范场与标量场的某些 BPS 方程之解的模空间。而这些 BPS 方程正是 $C$ 上的 Hitchin 方程。这两种紧化对应不同的极限，但我们预期：在交换 $C$ 与 $S^1$ 的相对长度标度时，低能物理不会发生相变。理由是：我们所研究的受 BPS 保护的量对 $C$ 上度量的共形标度不敏感——这要归功于下面要讲的拓扑扭转。因此，我们可以把三维 sigma 模型的靶空间 $\CM$ 等同于 $C$ 上 Hitchin 方程之解的模空间。这一点我们在第 \ref{sec:hiteq} 节中说明。

本文余下部分对 $\CM$ 所作的分析，只有在 $C$ 上带有某些缺陷算符（插入在点 $\CP_i$ 处）时才适用。关于这些缺陷，我们在第 \ref{sec:defects}–\ref{sec:bc-intersection} 节中作一些评论。

\subsubsection{$\IR^{1,5}$ 上的 $(2,0)$ 理论及其库仑分支}\label{subsubsec:20AndCoulomb}

我们从六维时空中的 $(2,0)$ 超共形理论讲起\cite{Witten:1995zh,Strominger:1996ac,Witten:1995em,Seiberg:1996vs,Seiberg:1997ax,Seiberg:1997zk}。这些理论具有 $osp(6,2\vert 4)$ 超共形不变性。模去拓扑上的微妙之处，它们由两类基本构件的乘积得到：一类是带 ADE 分类的相互作用理论，另一类是带阿贝尔「规范群」的自由理论。本文的主要兴趣在 A 系列的理论。从 M 理论的观点看，「规范群」为 $U(K)$ 的理论可以描述为 $K$ 个重合 M5-膜所成系统的解耦极限\cite{Strominger:1996ac}。我们经常把这一图像当作理解该理论性质的便捷捷径。

让我们简要回顾 $(2,0)$ 理论手征算符的一些信息；更多细节见 \cite{Aharony:1998an}，另可参看 \cite{Bhattacharya:2008zy}，尤其是第 31 页的表。存在一组算符基，它们在 $osp(6,2\vert 4)$ 的短表示中变换，并以 ADE 群 $\fg$ 的卡西米尔算符来标记。把这些卡西米尔算符记为 $k=1, \dots, r$。在第 $k$ 个短多重态中，我们将关注共形权最低的算符所张成的子空间 $V_k$。$V_k$ 是 $so(5)$ R 对称性的一个不可约表示，其共形权是 $\fg$ 的指数 $d_k$ 的两倍。

该理论有一个「库仑分支」，由这些手征算符的真空期望值参数化。在 $A_{K-1}$ 理论中，这个分支特别容易理解：它就是 $(\IR^5)^K / S_K$，参数化 $K$ 个 M5-膜在横向 $\IR^5$ 中彼此分离的组态。

在库仑分支上，理论含有 BPS 弦；几何上它们被描述为在分离的 M5-膜之间伸展的 M2-膜的边界。见图 \ref{fig:string}。

\insfig{4in}{string}{左：三个分离的 M5-膜，其中包括在两膜之间伸展的两段 M2-膜。右：$K=3$ 时 $A_{K-1}$ $(2,0)$ 理论中的对应图像；两段 M2-膜已被投影成弦段。}

把从膜 $i$ 与膜 $j$ 之间伸展的 M2-膜所得的弦称为 $ij$-弦。这些弦是有向的；反向则把 $ij$-弦与 $ji$-弦互换。BPS 条件要求这些弦是 $\IR^{5,1}$ 中的直线。BPS $ij$-弦的张力可由 M2-膜图像算出\footnote{M2-膜的张力是 $2\pi/\ell^3$，电磁对偶的 M5-膜的张力是 $2\pi/\ell^6$，其中 $\ell$ 是 11 维普朗克长度。}
\begin{equation}
T_{ij} = \frac{2 \pi}{\ell^3} |x_i-x_j|
\end{equation}
其中 $|x_i-x_j|$ 是 $\IR^5$ 中的距离（量纲为长度）。由 BPS 非重整化定理，人们预期这一几何图像其实给出精确的张力——即使在（本文最关心的）膜间距为 $\ell$ 量级的区域也是如此。

可以有超对称的「汇点」，在那里 $ij$、$jk$、$ki$ 型的三条弦相交。这种汇点在 $(2,0)$ 理论中的出现已在 \cite{Lee:2006gqa} 中指出。单条 $ij$-弦所保持的超对称，既依赖于弦在 M5-膜世界体中的斜率，也依赖于 $x_i-x_j$ 在 $\IR^5$ 中的方向。要找四分之一 BPS 汇点的条件，我们可以仿照 $(p,q)$ 弦汇点的标准设置\cite{Sen:1997xi, Aharony:1997bh, Dasgupta:1997pu}：把汇点平面内三条弦的相对斜率，与向量 $(x_i-x_j)$、$(x_j-x_k)$、$(x_k-x_i)$ 在 $\IR^5$ 中由三点 $x_i, x_j, x_k$ 所定平面内的相对斜率对应起来。注意，这自动给出汇点处的力学平衡：因为三段弦的张力向量此时与 $(x_i-x_j)$、$(x_j-x_k)$、$(x_k-x_i)$ 线性相关，故其和为零。

\subsubsection{$\IR^{1,3} \times C$ 上的 $(2,0)$ 理论及其库仑分支}

为了在 $C$ 上紧化 $(2,0)$ 理论而不破坏四维超对称，我们需要考虑一种部分扭转。$osp(6,2\vert 4)$ 的超庞加莱子代数的玻色部分为 $so(5,1) \oplus usp(4)\cong so(5,1) \oplus so(5)$。$Spin(1,5)$ 与 $Spin(5)$ 的旋量表示都是 $\IH^2$，其中 $\IH$ 是四元数。因此庞加莱超荷在 $so(5,1) \oplus so(5)$ 的 $(\IC^4\otimes \IC^4)_{+}$ 中变换，其中下标 $+$ 表示辛 Majorana 实在性约束。在 $C$ 上紧化后，我们保持子代数 $so(3,1)\oplus so(2)_C \oplus so(3) \oplus so(2)_R$，超荷在其中按
\begin{equation}
\left( \left((2,1)_{\half} \oplus (1,2)_{-\half} \right) \otimes
\left(2_{\half} \oplus 2_{-\half} \right) \right)_+
\end{equation}
变换。扭转的做法是：把 $so(2)_C\oplus so(2)_R$ 的对角 $so(2)$ 与 $C$ 的和乐代数等同起来；于是剩下的超荷在 $so(3,1) \oplus so(3) \oplus so(2)_C'$ 下按
\begin{equation}\label{eq:supercharge-decomp}
(2,1;2)_{1} \oplus (2,1;2)_0 \oplus (1,2;2)_0 \oplus (1,2;2)_{-1}
\end{equation}
变换。我们可以引入相应的超荷基：
\begin{equation}
Q_z^{\alpha A} ; Q^{\alpha A}; \bar Q^{\dot \alpha A}; \bar Q_{\bar z}^{\dot \alpha A}
\end{equation}
其中 $\alpha, \dot \alpha, A$ 都取 $1,2$，$z$ 是局部坐标，使得 $dz$ 带 $so(2)_C'$ 荷 $+1$。由于 (\ref{eq:supercharge-decomp}) 中间两项的超荷 $Q^{\alpha A}; \bar Q^{\dot \alpha A}$ 在 $so(2)_C$ 下不带荷，它们是良定义的四维超荷，因此我们保持了四维 $\CN=2$ 超对称。对易子 $\left\{ \bar Q^{\dot \alpha A}, \bar Q_{\bar z}^{\dot \beta B}\right\} \sim \epsilon^{\dot \alpha \dot \beta} \epsilon^{AB}\p_{\bar z}$ 在下面考虑被 $\bar Q^{\dot \alpha A}$ 湮灭的算符的手征环时会很有用。

四维理论的模空间本质上由六维理论的库仑分支作量纲约化得到。更确切地说，取 $so(5)$ R 对称性的一个嘉当子代数 $so(2)_R \oplus so(2)$，并令 $\CO_k$ 为 $V_k$ 中权为 $(d_k, 0)$ 的算符。该算符在该多重态中具有最大的 $so(2)_R$ 荷，因此必被任何带正 $so(2)_R$ 荷的超荷（例如 $\bar Q^{\dot \alpha A}$）湮灭。把真空期望值赋给这些手征算符，就到达 $d=4$ 的库仑分支。

扭转之后，$\CO_k$ 是 $C$ 上丛 $K^{\otimes d_k}$ 的截面，其真空期望值 $\inprod{\CO_k}$ 尤其如此。由于 $\CO_k$ 被 $\bar Q^{\dot \alpha A}$ 湮灭，而 $\bar Q^{\dot \alpha A}$-恰当算符的真空期望值为零\cite{Witten:1988ze}，故 $\inprod{\CO_k}$ 必被 $\bar\partial$ 湮灭。这是对 $\inprod{\CO_k}$ 的唯一条件，于是参数化这些真空期望值的 $d=4$ 库仑分支就是
\begin{equation}
\CB = \bigoplus_{k=1}^r H^0(C, K^{\otimes d_k}).
\end{equation}

在 $A_{K-1}$ 情形，$\CB$ 有一个漂亮的几何解释。我们考虑的是 $K$ 个 M5-膜所成系统的低能极限，这些膜包在某个超凯勒四维流形 $Q$ 内的全纯闭链 $C$ 上。要到库仑分支，我们把诸膜分开，使它们包在 $Q$ 内的另一条闭链 $\Sigma$ 上。\footnote{也可以把这些膜在三个平坦的横向方向上分开；那对应走到 $d=4$ 的希格斯分支，本文不涉及。} 与平坦空间的情形不同，我们将考虑膜无法完全分开的情形，于是 $\Sigma$ 是 $Q$ 内的一个\emph{连通}除子。取 $Q$（作为全纯辛流形）的全纯达布坐标 $(x,z)$，就可以把 $C$ 的一个邻域等同于全纯余切丛 $T^*C$。$\CB$ 中一点的选择，恰好对应选取方程
\begin{equation} \label{eq:sigma-eq}
x^K + \sum_{k=2}^K u_k(z) x^{K-k} = 0
\end{equation}
的系数 $u_k \in H^0(C, K^{\otimes k})$，$k = 2, \dots, K$，而该方程定义了 $\Sigma \subset T^* C$。

我们把坐标归一化，使得全纯辛形式为
\begin{equation}
\Omega = \frac{\ell^3}{2 \pi^2} dx \wedge dz.
\end{equation}
投影映射 $T^* C \to C$ 把 $\Sigma$ 实现为 $C$ 的 $K$ 重覆盖。第 $i$ 叶与第 $j$ 叶之间的距离是 $C$ 上的一个 1-形式，我们称之为 $\lambda_{ij}$。

下面我们将看到，$\Sigma$ 应等同于约化理论的 Seiberg–Witten 曲线。典范一形式
\begin{equation}
\lambda = x\,dz,
\end{equation}
限制到 $\Sigma$ 上，将等同于 Seiberg–Witten 微分。

\subsubsection{库仑分支上的 BPS 弦} \label{sec:bps-strings}

我们对库仑分支的几何图像，是读取 BPS 弦性质的一条便捷途径。与平坦空间的情形一样，局部我们有以叶对 $(i,j)$ 标记的 BPS $ij$-弦，以及三条此类弦相交的 BPS 汇点。$ij$-弦的 BPS 张力是 $\frac{1}{\pi} \abs{\lambda_{ij}}$；与平坦空间不同，这个张力依赖于 $C$ 上的点。$C$ 上还有一些特殊点，即 $\lambda_{ij} = 0$ 的 $ij$-\emph{分支点}，也就是 $\Sigma$ 的第 $i$ 叶与第 $j$ 叶相合之处。$ij$-弦可以终止于 $ij$-分支点。导出这一事实的一条捷径是回想这些态在 M5-膜图像中的描述：$ij$-弦是一张 M2-膜，它被连接第 $i$ 叶与第 $j$ 叶的线段所叶化；当这些线段收缩到零长时，膜便光滑地终止。见图 \ref{fig:m2-branch}。

\insfig{4in}{m2-branch}{左：一张在 M5-膜两叶之间伸展的 M2-膜的局部。该 M2-膜被「vertical」线段所叶化，每条这样的线段都落在 $T^* C$ 的单个纤维中。在两叶相合的分支点处，这些竖直线段收缩为零长。右：该 M2-膜到 $C$ 上的投影是 $(2,0)$ 理论中的一条弦，它终止于该分支点。}

沿曲线 $c$ 伸展的一段 $ij$-弦，其中心荷为
\begin{equation}\label{eq:Z-fragment}
Z = \frac{1}{\pi} \int_c \lambda_{ij}.
\end{equation}
另一方面，同一段弦的质量不过是张力的积分，
\begin{equation} \label{eq:stringmass}
M = \frac{1}{\pi} \int_c \abs{\lambda_{ij}}.
\end{equation}
于是 BPS 界 $M \ge \abs{Z}$ 取等，当且仅当 $\lambda_{ij}$ 沿曲线 $c$ 处处有相同相位 $\vartheta$；也就是说，若以 $\partial_t$ 记 $c$ 的切向量，我们要求
\begin{equation}\label{eq:BPS-WKB-curve}
\lambda_{ij} \cdot \partial_t \in e^{i \vartheta} \IR_+.
\end{equation}
此时 $Z = e^{i \vartheta} M$。

对多弦汇点，相位 $\vartheta$ 必须对所有弦股相同。这个条件等价于汇点处的无力条件：确实，代表 $ij$-股对汇点所施之力的向量可以写成复数 $e^{i \vartheta} \bar{\lambda}_{ij}$；既然三个 $\lambda_{ij}$ 之和为零，力之和也为零。

最简单的 BPS 物体是张在两个 $ij$-分支点之间的 $ij$-弦。在 M5-膜图像中它由一张圆盘代表；见图 \ref{fig:bps-disc}。这样的 M2-膜没有模。其量子化给出四维中一个 BPS 超多重态。

\insfig{5in}{bps-disc}{左：一张包在圆盘 $D$ 上的 M2-膜，它张在承载 M5-膜的曲线 $\Sigma$ 的两叶之间。右：在投影 $T^* C \to C$ 下，该圆盘投影为 $(2,0)$ 理论中的一条弦，后者终止于分支点。}

类似地，我们可以考虑沿 $C$ 中一条闭合回路 $c$ 伸展的 $ij$-弦。在 M5-膜图像中，这样的弦由一张柱面状的 M2-膜代表；见图 \ref{fig:bps-annulus}。这样的 M2-膜有一个玻色模及相应的费米零模。其量子化给出一个 BPS 矢量多重态。

\insfig{5in}{bps-annulus}{左：一张包在柱面上的 M2-膜，它张在承载 M5-膜的曲线 $\Sigma$ 的两叶之间。右：在投影 $T^* C \to C$ 下，该圆盘投影为 $(2,0)$ 理论中的一条闭弦。}

若 $K=2$，这些情形穷尽了所有可能；对态的这种描述曾在 \cite{Klemm:1996bj,Brandhuber:1996xk} 中用来研究 BPS 谱。对 $K>2$，还应考虑弦 web。同样把每一股 $c$ 提升为 $c_i - c_j$，所有提升之和给出 $\Sigma$ 中的一个闭链，它给出 BPS 粒子的荷。玻色模的个数等于 $c$ 中圈的个数 $\ell$。确实，$\ell = 1+(j-b)/2$，其中 $b$ 是分支端点的个数，$j$ 是三弦汇点的个数。每一股都可以垂直于自身移动，只不过连到分支点的每一股失去这个自由度；此外，每个汇点只有两个自由度，因而对终止于其上的三股再加一个约束。弦股的总数是 $(b+3j)/2$，于是这种朴素的玻色模计数给出 $(j-b)/2$。稍加努力就能看出，汇点处各约束的一个整体组合是平凡的，因而多出一个模，总共为 $\ell$（见 \cite{Aharony:1997bh} 式 23）。由超对称，我们应当还有 $\ell$ 个费米模。把这些费米模计入，我们预期量子化应给出自旋为 $\half(\ell+1)$ 的多重态。\footnote{由于模空间带有边界，视具体的边界条件而定，也可能出现自旋更低的多重态。}

最后让我们考虑这些 BPS 态的中心荷。$C$ 上代表 BPS $ij$-弦的有向曲线 $c$ 可以提升为 $\Sigma$ 上的一对曲线 $c_i$、$c_j$：即取 $c$ 在第 $i$ 叶与第 $j$ 叶上的原像，其中 $c_i$ 与 $c$ 定向相同，$c_j$ 定向相反。于是 \eqref{eq:Z-fragment} 可改写为
\begin{equation} \label{eq:central-lambda}
Z = \frac{1}{\pi} \int_c \lambda_{ij} = \frac{1}{\pi} \left( \int_{c_i} \lambda - \int_{c_j} \lambda \right).
\end{equation}
对代表某个 BPS 态的任何 $ij$-弦组合，提升后各曲线之和是 $\Sigma$ 上的一个闭链 $\gamma$；它正是代表该 BPS 态的 M2-膜的边界。我们得到的是
\begin{equation} \label{eq:central-charge}
Z = \frac{1}{\pi} \oint_\gamma \lambda.
\end{equation}
这个公式在本文余下部分至关重要。

顺便一提，也可以把 \eqref{eq:central-charge} 理解为斯托克斯定理与下述事实的推论：BPS 条件要求 M2-膜在 $Q$ 中是特殊拉格朗日的，即限制到该 M2-膜上时，全纯辛形式 $\Omega = \frac{\ell^3}{2 \pi^2} d \lambda$ 等于体积形式乘以 $e^{i \vartheta}$。在场论极限中，我们考虑的是紧贴 $T^* C \subset Q$ 之零截面 $C$ 的 M5-膜；在此极限下，特殊拉格朗日 M2-膜变得「vertical」，并退化为弦 web（相关讨论见 \cite{Mikhailov:1997jv}）。

\subsubsection{电荷格} \label{subsubsec:ChargeLattices}

这里正好停下来考虑 $d=4$ 理论中全部荷（味荷与规范荷）所成的格 $\hat\Gamma$。上面我们说过，BPS 态的荷由一个 1-圈 $\gamma \subset \Sigma$ 决定，它是终止于 M5-膜上的 M2-膜的边界。于是一个人可能天真地把电荷格等同于 $H_1(\Sigma; \IZ)$。这稍微有点不对，原因有二。

第一，$H_1(\Sigma; \IZ)$ 中并非所有类都能承载 BPS 态。这在 $K=2$ 的情形最容易看出：既然所有 BPS 态都来自连接两叶的膜，它们在两个操作的复合下不变——交换两叶，以及反转定向。由此，所有 BPS 荷都落在 $H_1(\Sigma; \IZ)$ 中在这一复合操作下不变的子格里。

第二，$H_1(\Sigma; \IZ)$ 中某些类应被视为物理上等价，因为从一个走到另一个不花能量。当 $C$ 上某处插入的缺陷处存在未破缺的非阿贝尔规范对称性时，这在 $K>2$ 时会发生；此时覆盖的若干叶在该缺陷附近变得等同，特别地，绕该缺陷的一个圈可以在这些叶之间自由移动。所以要得到物理的电荷格，我们必须取商。假定没有偶然简并，可以把这一商操作地描述为模去 $Z$ 的核，即把具有相同周期 $Z_\gamma$ 的圈 $\gamma$ 等同起来。于是总的来说，$\hat\Gamma$ 是 $H_1(\Sigma; \IZ)$ 的一个子商（一个子格的商）。

规范荷之格 $\Gamma$ 描述起来更简单：它就是 $H_1(\bar\Sigma, \IZ)$，其中 $\bar\Sigma$ 是把 $\Sigma$ 的穿刺点填上所得的紧黎曼面。味格 $\Gamma_{\rm flav} \subset \hat\Gamma$ 是 $H_1(\Sigma;\IZ)$ 上相交配对 $\inprod{,}$ 的根。$\Gamma_{\rm flav}$ 中任何元素都可由绕穿刺点的小圈之线性组合代表。$\hat\Gamma$ 是 $\Gamma$ 被 $\Gamma_{\rm flav}$ 的扩张，
\begin{equation}\label{eq:Lattice-ExactSeq}
0 \to \Gamma_{\rm flav} \to \hat\Gamma \to \Gamma \to 0.
\end{equation}

当我们在库仑分支 $\CB$ 中移动时，曲线 $\Sigma$ 随之变化，特别地 $H_1(\Sigma; \IZ)$ 可能在 $\Sigma$ 退化的轨迹周围有单值性。因此 $\hat\Gamma$ 与 $\Gamma$ 都不是固定的格，而是随 $\CB$ 变化的格的局部系统。相形之下，$\Gamma_{\rm flav}$ 是单值不变的。

\subsubsection{动能项}\label{subsubsec:KineticTerms}

为完整起见，让我们稍详细地讨论把 $(2,0)$ 理论在 $C$ 上紧化所得的四维拉氏量，并验证 $\Sigma$ 确实是 Seiberg–Witten 曲线、$\lambda$ 确实是 Seiberg–Witten 微分。

我们先从库仑分支的 M5-膜图像出发，计算红外四维规范理论中矢量多重态标量的动能项。M5-膜作用量的标量部分就是 DBI 作用量，即 $2 \pi / \ell^6$ 乘以体积形式。四维理论的一个真空由某条固定的全纯曲线 $\Sigma_0 \subset Q$ 确定。我们考虑 M5-膜的这样的组态：在 $\IR^{3,1}$ 的无穷远处趋于 $\IR^{3,1} \times \Sigma_0$，但在 $\IR^{3,1}$ 内部可以涨落。这种涨落（至少是足够温和的那些，而这里我们需要的也只是这些）由映射
\begin{equation}
f: \IR^{3,1} \times \Sigma_0 \to Q
\end{equation}
给出。令 $y\in \IR^{3,1}$，并定义 $f_y:=f(y,\cdot): \Sigma_0 \to Q$。我们要求当 $y\to \infty$ 时 $f_y$ 趋于恒等映射。

在低能极限下，我们只关注势能为零的映射，即要求 $\Sigma_y := f_y(\Sigma_0)$ 对一切 $y$ 都是体积极小的闭链。等价地，要求 $\Sigma_y$ 对一切 $y$ 都在 $Q$ 中全纯。于是可以定义 $\hat f: y\to \Sigma_y$ 为映入「$\Sigma$ 的一切全纯形变」所成空间 $\CD$ 的映射，
\begin{equation}
\hat f: \IR^{3,1} \to \CD.
\end{equation}
现在我们要计算这种映射的动能。我们在 $\hat f$ 于 $\IR^{3,1}$ 上\emph{缓变}的近似下工作，从而可以截断到变分 $\partial_\mu \hat f$ 的二阶。这些变分是 $\hat f(y) \in \CD$ 处切空间的元素，而该切空间又可等同于法丛的全纯截面空间 $T_{\hat f(y)}(\CD) \cong H^0(N(\Sigma_y))$。把 DBI 作用量按导数展开到二阶（减去常数部分），得
\begin{equation} \label{eq:KK-volume}
S = \frac{\pi}{\ell^6} \int_{\IR^{1,3}} d^4x \, \eta^{\mu\nu}
\inprod{\partial_\mu \hat f, \partial_\nu \hat f},
\end{equation}
其中 $\inprod{,}$ 是 $N(\Sigma_y)$ 的截面空间上标准的厄米度量，由 $Q$ 上的度量诱导而来。我们关心使 $S$ 有限的涨落；为此我们要求对每个 $\mu$，$\partial_\mu \hat f$ 都是\emph{可归一化的}，
\begin{equation}
\norm{\partial_\mu \hat f}^2 < \infty.
\end{equation}
可归一化模张成 $T\CD$ 的一个子空间，并在 $\CD$ 上定义一个可积分布。令 $\CB$ 为过 $\Sigma_0$ 的那片叶。我们称 $\CB$ 为 $\Sigma_0$ 的「可归一化形变」空间。于是我们把注意力限制在映射
\begin{equation}
\hat f: \IR^{3,1} \to \CB.
\end{equation}
上。令 $r$ 记可归一化模的个数，即 $\CB$ 的复维数。

至此我们发现，四维理论的标量部分是一个映入 $\CB$ 的 sigma 模型。六维张量多重态还含有一个带自对偶场强的阿贝尔 2-形式规范势，该场的 Kaluza–Klein 约化给出四维阿贝尔规范场。独立规范场的个数等于 $\Sigma$ 上可归一化调和 1-形式空间维数的一半，等价地等于 $\Sigma$ 上可归一化全纯 1-形式空间的维数。事实上这也等于 $r$。为看出这一点，让我们简要考察 $\Sigma$ 附近 $Q$ 上的度量。令 $z$ 为 $\Sigma$ 上的局部坐标，则可取横向坐标 $x'$，使 $\Sigma$ 是轨迹 $x' = 0$，且 $\Omega = \frac{\ell^3}{2 \pi^2} dx' \wedge dz$，而凯勒形式形如
\begin{equation}
\omega \vert_\Sigma = i g(z, \bar{z}) d z \wedge d \bar{z} + i \frac{\ell^6}{4 \pi^4} g^{-1}(z, \bar{z}) Dx' \wedge D\bar{x}',
\end{equation}
其中 $D$ 是 $N(\Sigma)$ 中的协变导数。现在，由伴随公式，存在同构 $N(\Sigma) \cong T^{*1,0}\Sigma$，用与 $\Omega$ 缩并来描述最为方便。此外，$H^0(N(\Sigma))$ 与 $H^0(T^{*1,0}\Sigma)$ 上分别有厄米内积
\begin{equation}
\langle v_1 \frac{\p}{\p x'}, v_2 \frac{\p}{\p x'} \rangle:=
\int_{\Sigma} \left( v_1 \bar v_2 \frac{\ell^6}{4\pi^4}
g^{-1}(z,\bar z) \right) g(z, \bar z) \frac{i}{2} dz d\bar z
\end{equation}
与
\begin{equation}
\langle \omega_1, \omega_2 \rangle:=  \frac{i}{2} \int_{\Sigma}
\omega_1 \overline{\omega_2}.
\end{equation}
取定这些内积后，与 $\Omega$ 的缩并是等距：
\begin{equation} \label{eq:inprod-equiv}
\norm{v \partial_{x'}}^2 = \frac{i}{2}\int_{\Sigma} \left( g(z,
\bar{z}) dz \wedge d\bar{z} \right) \frac{\ell^6}{4 \pi^4} g^{-1}(z,
\bar{z}) \abs{v}^2 = \frac{i}{2} \frac{\ell^6}{4 \pi^4}
\int_{\Sigma} \abs{v dz}^2 = \frac{\ell^6}{4 \pi^4} \norm{v dz}^2.
\end{equation}
特别地，$v dz$ 可归一化当且仅当 $v \partial_{x'}$ 可归一化。

总之，我们得到了一个四维阿贝尔 $\CN=2$ 超对称规范理论，含 $r$ 个矢量多重态。（回忆阿贝尔矢量多重态含一个复标量与一个规范场。）

我们还没有用标准方式描述这个 $\CN=2$ 理论。为此，设 $\Sigma$ 紧贴 $C$，并像上面那样把 $C$ 的一个邻域等同于 $T^* C$。要得到可归一化全纯 1-形式空间的一组好基，方便的做法是记住：它们正是那些可全纯延拓到 $\bar\Sigma$ 上的形式。于是在 $\bar\Sigma$ 上取一组 $A$、$B$ 圈基 $\{A^I, B_I\}$，$I=1,\dots, r$，并令 $\{\alpha_I\}$ 记对偶于 $A$ 圈的可归一化全纯 1-形式，
\begin{equation}
\oint_{A^I} \alpha_J = \delta^I_J.
\end{equation}
于是还可以照常定义 $\bar\Sigma$ 的周期矩阵
\begin{equation}
\tau_{IJ} := \oint_{B_I} \alpha_J.
\end{equation}
$\Sigma$ 上可归一化全纯 1-形式上的度量，与 $\bar\Sigma$ 上全纯 1-形式上的度量相同：
\begin{equation} \label{eq:metric-1-forms}
\inprod{\alpha_I, \alpha_J} = (\im \tau)_{IJ}.
\end{equation}

再定义
\begin{equation}
a^I := \frac{1}{\pi} \oint_{A^I} \lambda.
\end{equation}
这些 $a^I$ 给出 $\CB$ 上的局部坐标。我们想用这些坐标描述标量作用量 \eqref{eq:KK-volume}。为此需要更显式地理解 $T\CB$ 与 $H^0(N(\Sigma))_{\rm normalizable}$ 之间的等距。我们断言：若 $v_J dz = \pi \alpha_J$，则 $\frac{\p}{\p a^J} \to v_J \frac{\p}{\p x'}$。要看出这一点，注意 $v \frac{\p}{\p x'}$ 的作用是把曲面平移，使得
\begin{equation}
\pi \left( v \frac{\p}{\p x'}\right) a^I = \int_{A^{I'}}  \lambda
-\int_{A^{I}}   \lambda= \int_{W^I}   d\lambda = \int_{A^I} \iota( v
\frac{\p}{\p x'})  d \lambda =   \int_{A^I} v dz
\end{equation}
其中 $W^I$ 是把 $A^I$ 沿向量场 $v \frac{\p}{\p x'}$ 推挤所得的曲面的一小块无穷小部分。

现在，结合等距 \eqref{eq:inprod-equiv}，我们有
\begin{equation}
\p_\mu \hat f \cong \p_\mu a^I \frac{\p}{\p a^I} \to \p_\mu a^I
\frac{\ell^3}{2\pi} \alpha_I
\end{equation}
从而
\begin{equation} S = \frac{1}{4\pi} \int_{\IR^{1,3}}
d^4x \, \eta^{\mu\nu} (\im \tau)_{IJ} \partial_\mu a^I \partial_\nu
\bar{a}^J.
\end{equation}
这正是阿贝尔 $\CN=2$ 理论中动能项的标准形式；在这里我们\emph{把 $\Sigma$ 认同为 Seiberg–Witten 曲线，把 $\lambda$ 认同为 Seiberg–Witten 微分}。

\subsubsection{在 $C \times S^1$ 上紧化 $(2,0)$ 理论：Hitchin 方程} \label{sec:hiteq}

现在我们准备好考虑把理论从 $d=4$ 沿 $S^1$ 进一步量纲约化到 $d=3$ 所得的理论。

如前所述，我们理解这个理论的办法是回到六维，考虑\emph{先}在 $S^1$ 上、\emph{再}在 $C$ 上紧化。于是我们先在半径为 $R$ 的 $S^1$ 上紧化 $(2,0)$ 理论。所得理论在低能下是五维超对称 Yang–Mills 理论。我们感兴趣的是把这个五维理论在 $C$ 上进一步紧化，并作适当的拓扑扭转。

所得三维理论的模空间 $\CM$ 正是五维理论的 BPS 组态空间——这些组态还在 $\IR^3$ 中庞加莱不变。这些组态是什么？把超 Yang–Mills 理论的伴随标量记为 $Y^I$，$I=1, \dots, 5$，使得
\begin{equation} \label{eq:phiy}
\varphi := \half(Y^1 + i Y^2)
\end{equation}
带 $so(2)_R$ 荷 $+1$，而 $Y^{3,4,5}$ 荷为零。在扭转后的理论中，$\varphi = \varphi_z dz$ 是 $C$ 上的 $(1,0)$ 形式。于是 BPS 方程就是 $C$ 上规范场 $A = A_z dz + A_{\bar z} d\bar z$ 与伴随标量 $\varphi$ 的 Hitchin 方程：\footnote{证明 BPS 方程就是 Hitchin 方程的一条简洁途径如下。我们研究的是 $S^1 \times C \times \IR^3$ 上的一叠 M5-膜，等价地是 $C \times \IR^3$ 上的一叠 D4-膜，并寻找在 $\IR^3$ 平移下不变的 BPS 组态。众所周知，一叠 D4-膜上沿某个平移方向不变的 BPS 组态由自对偶 Yang–Mills 方程的解给出。而自对偶 Yang–Mills 方程在沿两个方向平移不变的组态上求值时，就变成 Hitchin 方程\cite{MR89a:32021}。}
\begin{equation} \label{eq:hitchin-3d-1}
F + R^2 [\varphi , \bar\varphi] = 0,
\end{equation}
\begin{equation} \label{eq:hitchin-3d-2}
\bar\p_A \varphi := d\bar z\left( \p_{\bar z} \varphi  + [ A_{\bar z}  ,\varphi]\right) = 0,
\end{equation}
\begin{equation} \label{eq:hitchin-3d-3}
\p_A \bar\varphi := dz\left( \p_z \bar\varphi  + [ A_z,\bar\varphi] \right) = 0.
\end{equation}

我们稍作插叙，解释 \eqref{eq:hitchin-3d-1} 中出现的、略显非常规的因子 $R^2$ 的来源。首先注意，在半径为 $R$ 的圆上约化 $(2,0)$ 理论后，得到的五维 SYM 作用量形如\footnote{由 $\varphi$ 与 $Y$ 之间的关系 \eqref{eq:phiy}、以及 $Z$ 与 $\varphi$ 之间的关系 \eqref{eq:central-charge} 可知：当 $Y^I = y^I \sigma^3$ 时，张在两个 M5-膜之间的开膜产生张力为 $\frac{1}{2\pi} y^I$ 的弦。既然膜的张力是 $2\pi/\ell^3$，可知沿横向到膜的实际距离是 $\frac{\ell^3}{4\pi^2} y^I$。检验 \eqref{eq:D4-action} 两项相对归一化的一个好办法，是计算 W 玻色子的质量，即两片错位的 D4-膜之间运行的弦的质量，结果为 $Ry^I$。而整体归一化可由约化六维 DBI 作用量得到。}
\begin{equation}\label{eq:D4-action}
S =  \frac{ R}{8\pi^2} \int_{\IR^{1,2}\times C} {\rm Tr} \left(
\frac{1}{R^2} F \wedge \star F + DY^I \wedge \star D Y^I + \cdots
\right)
\end{equation}
我们可以作重标度 $Y^I = R^{-1} \hat Y^I$ 以达到标准归一化：
\begin{equation}\label{eq:5deeAct}
S = \frac{1}{8\pi^2 R} \int_{\IR^{1,2}\times C} {\rm Tr} \left( F
\wedge \star F + D\hat Y^I \wedge \star D \hat Y^I + \cdots \right)
\end{equation}
正是在这个标架下，BPS 方程取与 $R$ 无关的形式 $F + [\hat \varphi, \overline{\hat \varphi}]\,dz\,d \bar z=0$。这就解释了 \eqref{eq:hitchin-3d-1} 中的因子 $R^2$。我们宁愿用 $\varphi$ 而不用重标度后的 $\hat \varphi$，理由是：用 $\varphi$ 表达时，希格斯场在奇点处的边界条件与 $R$ 无关——这一点后面会看到。

于是我们发现，$d=3$ 理论的模空间 $\CM$ 就是 $C$ 上 Hitchin 方程的解空间。在 \cite{Cherkis:1997aa,Gukov:1998kt,Kapustin:1998xn,Cherkis:2000cj} 中，人们用对偶关系论证：$SU(K)$ 的 $d=4$、$\CN=2$ 超对称场论紧化在圆上之后的模空间，应等同于某个特定的 Hitchin 系统。本节把他们的结果推广了。\footnote{我们感谢 Edward Witten 建议一般的 Hitchin 系统可以这样实现。} M5-膜与 Hitchin 系统之间的一般关系，Witten 已在 \cite{Witten:1997sc} 第 2.3 节中粗略勾勒过。

人们已经注意到，许多 $\CN=2$ 理论的 Seiberg–Witten 解可以用已知的复可积系统来理解，例如见 \cite{Gorsky:1995zq,Martinec:1995by,Donagi:1996cf,Itoyama:1995nv} 的讨论。这里我们构造了一大类一般理论，其相应的复可积系统是 Hitchin 系统。

关于 $\CM$ 的一般结构，我们再讲几句。基于一般的场论理由\cite{Seiberg:1996nz}，我们预期 $\CM$ 是超凯勒的，并且是 $d=4$ 理论库仑分支 $\CB$ 上的纤维化，一般纤维是紧环面。此外还预期这个环面纤维相对于 $\CM$ 上某个「特异」复结构是全纯子流形。在我们这里如何看出这一结构？

从 $\CM$ 到 $\CB$ 的投影容易描述：$\CB$ 由真空期望值 $\inprod{\CO_k}$ 参数化，而这些在五维 Yang–Mills 理论中被认同为 $\varphi$ 的独立卡西米尔不变量，所以投影不过是
\begin{equation}
(A, \varphi) \mapsto \{ \text{$\varphi$ 的卡西米尔不变量} \}.
\end{equation}
这个映射在数学文献中称为「Hitchin 纤维化」\cite{MR89a:32021}。它在一般点 $u \in \CB$ 上的纤维确实是一个阿贝尔簇，即投影 $\bar\Sigma_u \to C$ 的 Prym 簇，定义为雅可比簇之间相应映射 $J(\bar \Sigma_u) \to J(C)$ 的核。在重要的特例 $C = \IC\IP^1$ 中，$J(C)$ 平凡，于是它就简单地是 $J(\bar\Sigma_u)$。

在 $A_{K-1}$ 理论的情形，如前所述，$\inprod{\CO_k}$ 确定 Seiberg–Witten 曲线 $\Sigma \subset T^* C$。既然我们已把这些认同为卡西米尔不变量 $\Tr \varphi^k$，由 \eqref{eq:sigma-eq} 给出的 $\Sigma$ 正是由 $\varphi$ 决定的\emph{谱曲线}，
\begin{equation}\label{eq:speccurve}
\det (x\,dz - \varphi) = 0.
\end{equation}
换言之，$\Sigma$ 各叶在余切方向上的位置 $x_i$（$i = 1, \dots, K$）可以解释为矩阵值 1-形式场 $\varphi$ 的各个\emph{本征值}。（于是 \eqref{eq:sigma-eq} 中的系数 $u_k(z)$ 是本征值的初等对称函数，可以写成 $\langle \CO_k \rangle$ 的多项式。）

\subsubsection{缺陷与几何 Langlands} \label{sec:defects}

到目前为止，我们对 $C$ 究竟该是什么、以及该给场加什么边界条件，说得还有些含糊。\emph{A priori} 人们会认为最简单的可能是：就取 $C$ 为紧黎曼面，并要求场处处正则。然而事实证明，所得 $\CN=2$ 理论的 BPS 谱很难用本文的方法分析。不过有一个简单的修改能使事情简化：我们让 $C$ 是紧黎曼面，并在有限多个点上插入\emph{缺陷}。为了对允许何种缺陷有所认识，我们现在从三维短暂绕道到二维。

若把三维 sigma 模型进一步紧化在圆 $\tilde S^1$ 上，我们得到二维 $\CN = (4,4)$ 非线性 sigma 模型，靶空间仍为 $\CM$。这一进一步紧化与 Witten 等人关于几何 Langlands 纲领的工作\cite{geom-lang-n4,Gukov:2006jk,Frenkel:2007tx}接上了。

确实，总的来说我们把 $(2,0)$ 理论先在 $C$ 上、然后在 $S^1 \times \tilde S^1$ 上紧化了。假设我们以相反的顺序来做。把 ADE $(2,0)$ 理论在 $S^1 \times \tilde S^1$ 上约化，得到四维 $\CN=4$ 超对称的 ADE 规范理论。如 \cite{Harvey:1995tg,Bershadsky:1995vm} 所指出的，把 $\CN=4$ 理论在黎曼面 $C$ 上约化，导致一个映入 $\CM$ 的 sigma 模型。这个理论正是 \cite{geom-lang-n4,Gukov:2006jk,Frenkel:2007tx} 的出发点。

在几何 Langlands 纲领的数学文献中，有一类被深入研究过的缺陷（或「分歧」）可以在 $C$ 上引入。这些缺陷通过指定 Hitchin 方程之解的某些奇性来定义。物理上，它们可以解释为来自四维超 Yang–Mills 理论中的面算符。这类算符已在 \cite{Gukov:2006jk} 中研究过。在场论内部，它们可以用两种方式定义：或者规定路径积分中场具有特定的奇异行为，或者把某个二维 sigma 模型耦合到四维规范理论。另外，某些初等缺陷算符可以通过取相交膜组态的场论极限来定义：让给出四维规范理论的那叠 D3-膜与一些额外的 D3-膜沿一个余维 2 的轨迹相交。

归根到底，四维规范理论中的这些面算符应当来自 $(2,0)$ 理论中余维 2 的缺陷。然而，我们用来在四维定义缺陷的方法在六维不适用：我们没有六维 $(2,0)$ 理论的路径积分定义，也不知道如何把它与生活在余维 2 轨迹上的四维理论耦合起来。作为替代，我们主要通过缺陷在理论受保护算符中诱导的奇性来刻画它们。对本文而言这足够了，因为我们主要关心与 $\CN=2$ 场论相伴的 Hitchin 系统，而后者只依赖于受保护算符中的奇性。

我们将假定：\cite{Gukov:2006jk,Witten:2007td} 中用来产生带分歧的 Hitchin 系统的所有缺陷，都降自六维 $(2,0)$ 理论中的缺陷。对此假定我们可以提供一些部分证据，如下。在第 \ref{sec:bc-intersection} 节中，我们将把某些简单缺陷描述为特定 M5-膜交叠的场论极限。类似的、在受保护算符中诱导同类型奇性的缺陷，会以非紧 $C$ 上「无穷远处」边界条件的形式出现。有了这些简单缺陷，Hitchin 系统中任意复杂的分歧都可以由它们的碰撞以及其它极限程序产生，而这些在四维 $\CN=2$ 场论中似乎有简单的物理意义。正文中我们将考虑几个例子。

虽然很难给出缺陷算符的精确六维定义，但在 $S^1$ 上紧化后，至少应能给出五维的定义。确实，在红外，五维规范耦合趋于零，剩下的应是生活在缺陷处的一些 $\CN=4$ 三维 SCFT，它们与五维规范理论弱耦合。把它们的身份认出来将会很有意思。

\subsubsection{五膜交叠给出的边界条件} \label{sec:bc-intersection}

现在让我们简要讨论一种相对容易理解的特定缺陷。（下面第 3.2 节% TODO 合稿改 \ref{subsec:WittenConstruction}
中我们还会遇到一些更复杂的缺陷。）

我们再次考虑 $A_{K-1}$ 理论这一特例，它是用包在 $C$ 上的 $K$ 个 M5-膜实现的。现在考虑 $C$ 与另一条承载单个 M5-膜的曲线 $C_i$ 之间的一个简单横截相交。设两者在点 $\CP_i \in C$ 相交。于是在 $C$ 附近，$C_i$ 就是 $T^* C$ 在 $\CP_i$ 上的纤维。在库仑分支的一般点上，理论有一个（不可归一化的）形变能把 $C_i$ 与 $C$ 之间的相交抹平。具体说，取局部坐标 $z$ 使 $\CP_i$ 位于 $z = 0$，我们要求 $\Sigma$ 的某一叶在 $\CP_i$ 附近形如 $xz = \rho_i$，其中 $\rho_i \neq 0$。更不变量地说，回忆 $x$ 与 $\varphi$ 的本征值的对应，我们要求 $\varphi$ 在 $\CP_i$ 附近规范等价于
\begin{equation}\label{eq:commonsing}
\varphi = \frac{ 1 }{z} \begin{pmatrix} \rho_i & & & & & \\ & 0 & & & & \\ & & & \ddots & & \\ & & & & & 0 \end{pmatrix}  + \cdots
\end{equation}
我们还必须为规范场选取边界条件；一个与 Hitchin 方程相容的自然选择是
\begin{equation}\label{eq:Aboundryc}
A =\frac{1}{2i}  \left(\frac{d z}{z} -
\frac{d \bar z}{\bar z}\right) \begin{pmatrix} \alpha_i & & & & & \\ & 0 & & & & \\ & & & \ddots & & \\ & & & & & 0 \end{pmatrix} + \cdots
\end{equation}
其中 $\alpha_i$ 是纯虚数。

正如我们后面在例子中会看到的，$\rho_i$ 有一个自然的解释：它是四维规范理论中某个 $U(1)$ 味对称性的质量参数。此外，在约化到三维之后，$\rho_i, \alpha_i$ 中编码的三个实参数可以解释为某个 $U(1)$ 味对称性矢量多重态中三个标量场的真空期望值。

方程 (\ref{eq:commonsing}) 在极限 $\rho_i \to 0$ 下应小心解释；下面第 3.2.6 节% TODO 合稿改 \ref{sec:couplings}
中我们会进一步讨论这一点。

\begin{yizhu}
原文中 $\ell$ 一符两用：既是 11 维普朗克长度（式 (3.1) 的脚注），又是弦 web 的圈数（本节 3.1.3 的 $\ell = 1+(j-b)/2$）。译文一仍其旧，请读者按上下文区分。
\end{yizhu}
